\icode{N15-MC-SI-L3-261026-002}
\nceeID[6]{2015新课标\Rmnum{2}}{16}
由于卫星的发射场不在赤道上，同步卫星发射后需要从转移轨道经过调整再进入地球同步轨道。当卫星在转移轨道上飞经赤道上空时，发动机点火，给卫星一附加速度，使卫星沿同步轨道运行。已知同步卫星的环绕速度约为\qty{3.1e3}{m/s}，某次发射卫星飞经赤道上空时的速度为\qty{1.55e3}{m/s}，此时卫星的高度与同步轨道的高度相同，转移轨道和同步轨道的夹角为\qty{30}{\degree}，如图所示，发动机给卫星的附加速度的方向和大小约为\paren[]

\begin{tikzpicture}[font=\scriptsize]
	\draw[dashed,thin,cyan] (0,1.5)node[above]{ \(N\)}--(0,-1.5)node[below]{ \(S\)} (-2,0)node[above]{ \(W\)}node[left]{ 赤道}--(2,0)node[above]{ \(E\)}node[right]{ 同步轨道} (30:1.5)node[above]{ 转移轨道}--(-150:1.5);
	\draw[->,magenta] (0,0)--(1.55,0);
	\draw[->,blue] (0,0)--(30:.775);
	\draw (0,0) circle(1);
	\draw[fill=gray!50] (0,0) circle(.1);
\end{tikzpicture}
\begin{choices}
	\item 西偏北方向，\qty{1.9e3}{m/s}
	\item 东偏南方向，\qty{1.9e3}{m/s}
	\item 西偏北方向，\qty{2.7e3}{m/s}
	\item 东偏南方向，\qty{2.7e3}{m/s}
\end{choices}
\begin{answer}
	B\qquad 提示：已知合速度为同步卫星的线速度，一个分速度是在转移轨道上的速度，另一个分速度待求，运用速度合成的平行四边形法则求解即可
\end{answer}